% MG05a. After the extension/provenance paragraph following
% act21:tpk:logaritmo; before the reading-and-motion order paragraph.
\subsubsection{Gauges of the multiplicative observable and generator normalization}
\label{mg05:calibres}

The choice of discrete exponent admits a finite comparison in which
the APP support, the initial conditions of both cursors, the tritic
selector, the reversal calendar, and the emission rule remain fixed.
The variation is confined to the observable on the multiplicative
sheet. Let \(\ell_{2}\) and \(\ell_{5}\) be the logarithms of
\((\mathbb Z/9\mathbb Z)^\times\) with generators \(2\) and
\(5\), respectively, extended by zero on \(\{3,6,9\}\).
They are compared with \(v_3\), evaluated on the positive
representatives \(1,\ldots,9\), and with the identity reading on
those same representatives.

\begin{proposition}[Normalization of the multiplicative logarithm]
\label{mg05:normalizacion}
The two isomorphisms from the unit group modulo nine to
\(\mathbb Z/6\mathbb Z\), normalized by sending a generator to
\(1\), are \(\ell_2\) and \(\ell_5\). They satisfy
\[
 \ell_5(x)\equiv-\ell_2(x)\pmod6
 \qquad\bigl(x\in(\mathbb Z/9\mathbb Z)^\times\bigr).
\]
Choosing the least positive generator fixes \(\ell_2=\ell_9\),
with the values in \eqref{act21:tpk:logaritmo}.
\end{proposition}
\begin{proof}
The powers of \(2\) traverse \(1,2,4,8,7,5\) and close at the
sixth step. The generators of a cyclic group of order six correspond
to exponents coprime to six: \(1\) and \(5\). Their representatives
are \(2\) and \(2^5\equiv5\pmod9\).
Since \(5\equiv2^{-1}\pmod9\), expressing a unit as a power of
\(5\) changes the sign of its exponent in base \(2\).
Normalization by the least positive representative selects \(2\).
\end{proof}

The homomorphism condition alone also admits maps with proper image.
The proposition therefore distinguishes a faithful representation of
the unit group from its quotients. Extension by zero also specifies
the treatment of the three noninvertible residues; the provenance
record retains which residue was observed.

The comparison of the four observables traverses the same
\(104\,976\) initial conditions and groups their emissions over six
windows. Writing \(E_g\) for the emitter with multiplicative
observable \(g\), one obtains the following image census.
\begin{table}[htbp]
\centering
\small
\begin{tabular}{@{}lrrrl@{}}
\toprule
Observable & \(|\operatorname{im}E_g|\) & Zero roots
 & Additive law & Normalization\\
\midrule
\(\ell_2\)&468&yes&yes&least generator\\
\(\ell_5\)&468&yes&yes&inverse generator\\
\(v_3\)&144&no&yes&zero unit image\\
\(\operatorname{id}\)&684&no&no&positive residue\\
\bottomrule
\end{tabular}
\caption{Comparison of four gauges on the same transport and the same domain of initial conditions.}
\label{mg05:tabla-calibres}
\end{table}

Both logarithms satisfy the additive exponent law. The valuation
\(v_3\) vanishes on every unit but takes the values \(1,1,2\) on
\(3,6,9\); this accounts for the difference between its two control
columns. The identity retains the positive residue and fails the
homomorphism law into \(\mathbb Z/6\mathbb Z\); for example, at
the unit \(1\), the required equality would be \(1\equiv2\pmod6\).
The census of \(468\) emissions is shared by both logarithms and
therefore retains the inversion ambiguity. The least-generator
criterion resolves that ambiguity before any regional reading of
constants. The cardinalities in the table describe the exhaustive
evaluation of these four observables on the fixed domain; their
role is to test the emitter's operational normalization.
